\documentclass[10pt,a4paper]{amsart}
\usepackage[margin=19mm]{geometry}
\usepackage{amsmath,amssymb,mathtools}
\usepackage[scr=rsfs,cal=euler]{mathalfa}
\usepackage{xfrac}
\usepackage[unicode,hidelinks,bookmarks=false]{hyperref}
\newcommand{\ie}{i.e.\ }
\newcommand{\cf}{cf.\ }
\newcommand{\resp}{respectively\ }
\newcommand{\Subsection}{\subsection}
\providecommand{\nobreakdash}{\nobreak\hbox{-}}
\setlength{\emergencystretch}{2em}
\multlinegap0pt
\usepackage{amssymb}
\usepackage[shortlabels]{enumitem}
\newtheorem{proposition}{Proposition}
\newcommand{\N}{\mathbb N}
\newcommand{\R}{\mathbb R}
\title{Metric-geometric Chow theorem}
\author{Jos\'e Edson Sampaio}
\date{Source: arXiv:2608.27810v1 (CC BY 4.0)}
\begin{document}
\maketitle

\noindent\emph{Source and licence.} Metric-geometric Chow theorem. Version arXiv:2608.27810v1 (CC BY 4.0), distributed by arXiv under the Creative Commons Attribution 4.0 International licence, http://creativecommons.org/licenses/by/4.0/, as recorded in the arXiv licence metadata for this exact version and shown on the abstract page https://arxiv.org/abs/2608.27810v1. Full licence notice: \texttt{licenses/arxiv-2608.27810-CC-BY-4.0.txt}.\par\smallskip

\noindent\textit{Editorial scope.} Counted unit: the article's proposition selection\_lemma (source0.tex lines 330--336). The article's complete original proof of it is delivered with it (source0.tex lines 337--375, one \texttt{proof} environment of 39 lines). No mathematics is written, repaired or re-derived: the statement and the proof are the article's own lines, in the article's own notation, and no retained line is substituted or re-flowed.  No further source-local context is needed: every cross-reference of every retained line resolves inside the two delivered spans.  Nothing was omitted from any retained span, so the package carries no omission record.  No retained line contains a citation, so the wrapper carries no bibliography and no bibliography entry.  No retained line contains a figure environment, an image-inclusion command or drawing code, so no image is delivered and the package declares no asset dependency.  Editorial formatting only, in addition: an AMS \texttt{amsart} preamble that restates the article's own declarations and macros used by the retained lines, namely the environments \texttt{proposition} on separate counters, together with the standard packages those lines need.  The retained environments print in this document as Proposition 1 in order of appearance, whereas the article prints the same items with the numbering of its own full document.  The complete native source of the article as distributed in the arXiv e-print for this exact version is bundled unchanged as \texttt{sources/208742/source0.tex}.
\begin{proposition}\label{selection_lemma}
Let $Z\subset \R^n$ be an unbounded definable set in an o-minimal structure on $\R$. Then $C(Z,\infty)$ is definable, $\dim C(Z,\infty)\leq \dim Z$, and 
$$
C(Z,\infty)=\{w\in \R^n; \exists\ \gamma\in C^0((\varepsilon ,+\infty ), Z) \,\,  s.t.\,\,\gamma(t)=tw+o_{\infty }(t)\},
$$ 
where $g(t)=o_{\infty }(t)$ means $\lim\limits _{t\to +\infty }\frac{g(t)}{t}=0$ and $C^0((\varepsilon ,+\infty ), Z)$ is the set of all continuous functions from $(\varepsilon ,+\infty )$ to $Z$. 
\end{proposition}

\begin{proof}
Suppose that $w\in\R^n$ is a tangent vector of $Z$ at infinity. Let us consider the semialgebraic mapping $\iota\colon\R^n\setminus\{0\}\to \R^n\setminus\{0\}$ given by $\iota(x)=\frac{x}{\|x\|^2}$ and denote $X=\iota(Z\setminus \{0\})$.  Since $Z$ is an unbounded set, the origin is a non-isolated point of $\overline{X}$. Let $\rho\colon \mathbb{S}^{n-1}\times [0,+\infty )\to \R^n\setminus \{0\}$ be the mapping given by $\rho(x,t)=tx$. We see that $\rho|_{\mathbb{S}^{n-1}\times (0,+\infty )}:\mathbb{S}^{n-1}\times (0,+\infty )\to \R^n\setminus \{0\}$ is a semialgebraic homeomorphism with inverse mapping $\rho^{-1}\colon\R^n\setminus \{0\}\to \mathbb{S}^{n-1}\times (0,+\infty )$ given by $\rho^{-1}(x)=(\frac{x}{\|x\|},\|x\|)$. Therefore, the sets $Y=\rho^{-1}(X)\subset \mathbb{S}^{n-1}\times [0,+\infty )$ and $\overline{Y}$ are definable sets in an o-minimal structure on $\R$. We are going to consider two cases:

\bigskip

\noindent 1) Case $w\not=0$. Since $w$ is a tangent vector of $Z$ at infinity, there are sequences $\{s_k\}_{k\in \N}$ of positive real numbers and $\{z_k\}_{k\in \N}\subset Z$ such that $\lim\limits _{k\to +\infty }\|z_k\|=+\infty $ and $\lim\limits _{k\to +\infty }\frac{1}{s_k}z_k=w$. Thus, for each $k\in \N$, let us define $x_k=\iota(z_k)$. In this case,  $v:=\lim\limits _{k\to \infty } s_k\cdot x_k= \frac{w}{\|w\|^2}$. In particular, $\lim\limits _{k\to \infty } \frac{x_k}{\|x_k\|}=\frac{w}{\|w\|}=\frac{v}{\|v\|}$ and $u=(\frac{v}{\|v\|},0)\in \overline{Y}$. Then, by the Curve Selection Lemma, there exists a continuous curve $\beta:[0,\delta)\to \overline{Y}$ such that $\beta(0)=u$ and $\beta((0,\delta ))\subset Y$. By writing $\beta(t)=(x(t),s(t))$, we obtain that $s:[0,\delta )\to \R$ is a definable and non-constant function such that $s(0)=0$ and $s(t)>0$ if $t\in (0,\delta )$. By the Monotonicity Lemma, one can suppose that $s$ is $C^1$ and  strictly increasing in the domain $(0,\delta)$. Hence, $s\colon [0,\delta/2]\to [0,\delta' ]$ is a homeomorphism, where $\delta' =s(\frac{\delta}{2})$. Let us define $\alpha\colon [0,r)\to \overline{Y}$ by
$$
\alpha(t)=\rho\circ \beta\circ s^{-1}(t\|v\|)=\rho(x(s^{-1}(t\|v\|)),s(s^{-1}(t\|v\|)))=t\|v\|x(s^{-1}(t\|v\|)),
$$
where $r=\min\{\frac{\delta'}{\|v\|},\delta'\}$. Therefore,
$$
\lim\limits _{t\to 0^+}\frac{\alpha(t)}{t}=\lim\limits _{t\to 0^+}\frac{t\|v\|x(s^{-1}(t))}{t}=\lim\limits _{t\to 0^+}\|v\|x(s^{-1}(t))=\|v\|x(0)=v,
$$
and thus $\alpha(t)=tv+o(t)$. Finally, by defining $\gamma\colon(\frac{1}{r},+\infty )\to Z$ in the following way $\gamma(t)=\iota^{-1}(\alpha(\frac{1}{t}))$, we obtain
\begin{eqnarray*}
\gamma(t) &=& \frac{\frac{1}{t}v+o(\frac{1}{t})}{\|\frac{1}{t}v+o(\frac{1}{t})\|^2}=t\textstyle{\frac{v}{\|v\|^2}}+o_{\infty }(t)\\
		  &=& tw+o_{\infty }(t).
\end{eqnarray*}
Since $\gamma$ is a composition of continuous mappings,  $\gamma$ is a continuous mapping as well.

\bigskip

\noindent 2) Case $w=0$. In this case,  let $\{x_k\}_{k\in\N}\subset Z$ be a sequence such that  $\lim\limits _{k\to +\infty }\|x_k\|=+\infty $ (this sequence exists because $Z$ is unbounded). Thus, $\{\frac{x_k}{\|x_k\|}\}_{k\in\N}$ is, up to taking a subsequence, a convergent sequence. Let $v\in  \R^n$ be the limit of this sequence, i.e., $\lim\limits _{k\to \infty }\frac{x_k}{\|x_k\|}=v$. Likewise, as was done in Case 1, one can show that there exists a continuous curve $\gamma\colon (\varepsilon,+\infty )\to Z$ such that $\gamma(t)=tv+o_{\infty }(t)$. Let us define $\widetilde\gamma\colon (\varepsilon^{2}, +\infty )\to Z$ by $\widetilde{\gamma}(t)=\gamma(t^{\frac{1}{2}})$. Thus, we have $\widetilde{\gamma}(t)=o_{\infty }(t)=tw+o_{\infty }(t)$.

\bigskip

Reciprocally, if there exists a continuous curve $\gamma:(\varepsilon,+\infty )\to Z$ such that $\lim\limits _{t\to +\infty }|\gamma(t)|=+\infty$ and $\gamma(t)=tw+o_{\infty }(t)\}$, then for each $k\in \N$ we define $s_k=\varepsilon+k+1$ and $z_k=\gamma(s_k)$. Thus, it is clear that $w$ is a tangent vector of $Z$ at infinity, since $\lim\limits _{k\to +\infty }\|z_k\|=+\infty $ and $\lim\limits _{k\to +\infty }\frac{1}{s_k}z_k=w$.


Therefore, 
$$
C(Z,\infty)=\{w\in \R^n; \exists\ \gamma\in C^0((\varepsilon ,+\infty ), Z) \,\,  s.t.\,\,\gamma(t)=tw+o_{\infty }(t)\}.
$$ 
Moreover, we have also proved that $(C(Z,\infty)\cap \mathbb{S}^{n-1})\times \{0\}\subset \overline{Y}\setminus Y$. Then 
$$
\dim(C(Z,\infty)\cap \mathbb{S}^{n-1})\leq \dim(\overline{Y}\setminus Y)<\dim Y=\dim Z,
$$
which implies $\dim(C(Z,\infty))\leq \dim Z$.
\end{proof}

\end{document}
